\subsection{A restriction on arithmetic towers with noetherian integral closure}

\begin{proposition}[Proved: the Galois tower is virtually procyclic]
\label{C:A2-Galois-tower-procyclic}
Let $F$ be a number field, let $S$ be a finite set of finite primes
of $F$, and let $K/F$ be an algebraic Galois extension.  Let $R$ be
the integral closure of $\mathcal O_F[S^{-1}]$ in $K$.  If $R$ is
noetherian, then $\operatorname{Gal}(K/F)$ has an open procyclic
subgroup.  In particular, if $K/F$ is an infinite pro-$p$ extension,
there is a finite intermediate field $F'/F$ such that $K/F'$ is a
$\mathbf Z_p$-extension.
\end{proposition}

\begin{proof}
Write $A=\mathcal O_F[S^{-1}]$.  The ring $R$ is a normal domain,
integral over the Dedekind domain $A$, with fraction field $K$.
Its nonzero prime ideals are maximal.  Thus, under the stated
noetherianity hypothesis, $R$ is Dedekind.

Fix a nonzero prime $\mathfrak q$ of $A$.  Choose
$0\ne a\in\mathfrak q$.  Every prime of $R$ over $\mathfrak q$
contains $a$, so the factorization of $aR$ in the Dedekind domain
$R$ shows that there are only finitely many such primes.  Choose
one, $Q$, and let $D\subseteq G=\operatorname{Gal}(K/F)$ be its
decomposition subgroup.  The Galois action on the primes over
$\mathfrak q$ is transitive, by the finite-extension statement
and compactness.  Consequently $D$ is open in $G$.

We also record explicitly the ramification consequence of
noetherianity.  The local ring $R_Q$ is a discrete valuation ring;
normalize its valuation $v_Q$ to have value group $\mathbf Z$.
If $\pi\in F$ is a uniformizer at $\mathfrak q$, put
$e_\infty=v_Q(\pi)>0$.  For a finite intermediate field
$F\subseteq E\subseteq K$ and the induced prime $Q_E$, the
restriction of $v_Q$ to $E^\times$ is a positive integral
multiple of the normalized valuation at $Q_E$.  It follows that
\[
                 e(Q_E/\mathfrak q)\mid e_\infty.
\]
In particular, all finite-level ramification indices along $Q$
are bounded by the same integer.

For finite Galois intermediate fields, the inertia group has
order equal to that ramification index, and the quotient of the
decomposition group by inertia is the Galois group of the finite
residue-field extension.  We use the local decomposition and
inertia statements in \cite[Theorem~7.58 and
Propositions~8.10--8.11, pp.~131, 139--140]{C_MilneANT2020}.
Passing to inverse limits therefore gives an exact sequence
\[
  1\longrightarrow I\longrightarrow D\longrightarrow
  \operatorname{Gal}(\kappa(Q)/\kappa(\mathfrak q))
  \longrightarrow1,
\]
where $I$ is finite: all its finite quotient orders are bounded
by $e_\infty$.  The group on the right is procyclic because
$\kappa(\mathfrak q)$ is finite.

Choose a topological generator $\overline\sigma$ of $D/I$ and
a lift $\sigma\in D$.  The closed subgroup
$H=\overline{\langle\sigma\rangle}$ is procyclic.  Its image
in $D/I$ is compact and contains the dense cyclic subgroup
generated by $\overline\sigma$, so that image is all of $D/I$.
Hence $D=HI$ and $[D:H]\le |I|$.  A closed finite-index
subgroup of a profinite group is open, so $H$ is open in $G$.
If $G$ is infinite pro-$p$, then $H$ is an infinite procyclic
pro-$p$ group, hence isomorphic to $\mathbf Z_p$.  Its fixed
field is the required finite extension $F'$.
\end{proof}

\begin{remark}[Exact scope of the obstruction]
We did not assume that ramification is confined to $S$, or even
that there are only finitely many ramified primes.  One retained
prime already gives the conclusion.  The proposition concerns the
\emph{full integral closure} in a Galois tower.  A non-Galois
tower, or an overring obtained by discarding a non-Galois-invariant
collection of primes, is outside its conclusion.  Thus a
Golod--Shafarevich Galois tower with non-virtually-procyclic group
cannot directly furnish the sought noetherian Dedekind ring by
taking its full ring of $S$-integers.
\end{remark}

\begin{lemma}[Proved: the exact persistence requirement]
\label{C:A2-class-persistence}
Let $F_0\subset F_1\subset\cdots$ be an algebraic-number-field
tower, let $S_n$ be the primes above a fixed finite set $S_0$,
and put $A_n=\mathcal O_{F_n}[S_n^{-1}]$ and
$A_\infty=\varinjlim_n A_n$.  Then
\[
 \operatorname{Pic}(A_\infty)/p
   =\varinjlim_n\bigl(\operatorname{Pic}(A_n)/p\bigr),
\]
where the transition maps extend ideals.  A class
$c\in\operatorname{Pic}(A_n)$ survives nontrivially in this
quotient if and only if its extension to
$\operatorname{Pic}(A_m)$ belongs to no subgroup
$p\operatorname{Pic}(A_m)$, for every $m\ge n$.
More generally, a finite list is linearly independent in the
limit if and only if no nonzero $\mathbf F_p$-linear combination
of it becomes $p$-divisible at any later stage.
\end{lemma}

\begin{proof}
Invertible modules and their isomorphisms descend through a
filtered colimit of rings: we descend a finite projective
presentation, its rank-one condition, and an inverse module with
its evaluation isomorphism.  Therefore
$\operatorname{Pic}(A_\infty)=\varinjlim_n\operatorname{Pic}(A_n)$.
Taking the cokernel of multiplication by $p$ commutes with the
filtered colimit.  The stated survival and independence tests
are the defining zero test in a filtered colimit of vector spaces.
\end{proof}

If $A_\infty$ is Dedekind and this vector space is infinite
dimensional, deleting any finite set of closed points leaves it
nonzero: localization quotients its Picard group by the subgroup
generated by those finitely many prime classes.  By contrast,
unbounded finite-stage class-group ranks, or injectivity before
reduction modulo $p$, do not establish the persistence condition.

\paragraph{Former construction gate, now resolved.}
Successive non-Galois degree-$p$ steps could in principle be
chosen so that every fixed rational prime is eventually
unramified and inert at all subsequent steps.  The resulting
finite splitting and bounded ramification are exactly the
noetherianity conditions in Lemma~\ref{C:A2-tower-Dedekind};
growing local $p$-degrees are suited to killing Brauer classes.
For $p=2$ the quadratic-step construction in
\texttt{A2\_quadratic\_tower\_counterexample.tex} now combines these
local prescriptions with injectivity modulo two at every step and
doubles the class-group rank. Its full arithmetic and scheme-theoretic
proof supplies a counterexample to A.2. The Galois-tower obstruction
above remains valid; this constructed tower is not Galois over its
initial number field.
